\section{A filtered module on the boundary graph}
\label{bl:sec:hodge}

The tensor $v_k$ of Section~\ref{bl:cov:obstruction-contract} has
coefficients in $R^1h_*\omega_R(D_h)$. We construct a filtered right
$\mathcal D_B$-module whose lowest piece is this entire coherent
sheaf, so that differentiation in base coordinates gives a first
symbol for the tensor. We then prove a vanishing statement for the
kernel of that symbol on a projective base. Section~\ref{bl:sec:frames}
will apply it to the positive frame weight $a+k$.

Temporarily isolate the graph calculation from the covers. Let $Y$ be
a smooth complex quasi-projective variety, let $R+D_h$ be
a reduced simple normal crossing divisor with no common component
between $R$ and $D_h$, and let $h:R\to B$ be a projective morphism
to a smooth quasi-projective variety. Write $b=\dim B$, and let
$Q\subset Y\times B$ be its graph. The projection is not required
to be proper, but its restriction to $Q$ is projective:
\[
\begin{tikzcd}[column sep=large,row sep=large]
R \arrow[r,hook,"{(\iota,h)}"] \arrow[dr,"h"'] &
Y\times B \arrow[d,"\operatorname{pr}_B"]\\
& B.
\end{tikzcd}
\]
Throughout this section $D_h$ also denotes its pullback to the product.

\subsection{Conventions and the lowest filtration piece}

We use right $\mathcal D$-modules underlying graded-polarizable
algebraic mixed Hodge modules with rational structures. On a smooth
$d$-fold we twist the constant Hodge module by $d$, so that its
underlying canonical sheaf starts at $F_0$, rather than $F_{-d}$.
All subsequent localization and support functors carry this
normalization. The required functorial statements are the
localization and support operations, exterior products, smooth
inverse images, and projective direct-image strictness in
\cite{Saito1988,Saito1990,Saito2013}; conventions are also explained
in~\cite{Schnell2014,Popa2016}.

Define
\begin{equation}\label{bl:hod:modules}
 \mathcal M=\mathcal H^{b+1}_{[Q]}
       \bigl(\omega_{Y\times B}(*D_h)\bigr),
 \qquad
 \mathcal N=\mathcal H^1\operatorname{pr}_{B,+}\mathcal M.
\end{equation}
Here the brackets denote algebraic local cohomology, and the plus
sign denotes differential-module direct image. Localization and
cohomology with supports show that $\mathcal M$ is the underlying
filtered module of a mixed Hodge module. Its direct image has
projective support, as above.

For a right module, the relative filtered de Rham complex for the
projection has term
\begin{equation}\label{bl:hod:relative-dr}
 F_{j-i}\mathcal M\otimes\bigwedge^i T_Y
 \quad\text{in degree }-i.
\end{equation}
The tangent sheaf here is pulled back from $Y$. In particular its
last term, in degree zero, is $F_j\mathcal M$.

\begin{lemma}[Lowest piece and strict insertion]\label{bl:lem:lowest-piece}
With the normalization above,
\begin{align}
 F_{<0}\mathcal M&=0,&
 F_0\mathcal M&=(\iota,h)_*\omega_R(D_h),
 \label{bl:hod:lowest-source}\\
 F_{<0}\mathcal N&=0,&
 F_0\mathcal N&=R^1h_*\omega_R(D_h).
 \label{bl:hod:lowest-target}
\end{align}
The equality for $F_0\mathcal N$ identifies it with a subsheaf of $\mathcal N$.
It is induced by simple graph residues and the inclusion of the
last term of the relative de Rham complex. The latter inclusion
commutes with right differentiation in base coordinates.
\end{lemma}

\begin{proof}
Work locally along $Q$. Take simple normal crossing coordinates
for $R+D_h$, write $r_0$ for the reduced equation of $R$, and take
coordinates $t_1,\ldots,t_b$ on the base. The functions $h^*t_i$ on
$R$ lift locally to functions $v_i$ on $Y$. Thus $Q$ is defined by
\[
 r_0=0,\qquad \psi_i=t_i-v_i=0\quad(1\leq i\leq b).
\]
The $\psi_i$, together with the boundary coordinates, are independent.
The support sequence is a hypersurface followed by $b$ smooth
support coordinates, which explains the index $b+1$ in
\eqref{bl:hod:modules}.

Before adding the graph coordinates, the boundary module is
\begin{equation}\label{bl:hod:boundary-module}
 \omega_Y(*(R+D_h))/\omega_Y(*D_h).
\end{equation}
In one coordinate, localization of the normalized constant begins
with simple poles at $F_0$; $F_j$ permits $j$ further derivatives.
The quotient by regular forms is the residue delta module, whose
simple residue is again at level zero. This is the usual
localization--residue exact sequence with its strict Hodge filtration.
Taking exterior products gives the same rule for a normal crossing
divisor: one begins with simple poles in all the indicated
coordinates, and allows a total of at most $j$ extra denominator
exponents; see also~\cite[Section~3 and Proposition~8.2]{MustataPopa2019}.
Quotienting out the part regular in the $R$-support
direction gives~\eqref{bl:hod:boundary-module}. The $b$ independent
graph residues add free normal derivatives with the same order rule.

It follows that $F_{<0}\mathcal M=0$, and its lowest piece is
\[
 \omega_Y(R+D_h)/\omega_Y(D_h)=\omega_R(D_h),
\]
inserted along the graph by the class
\begin{equation}\label{bl:hod:simple-graph}
 \alpha\longmapsto
 \left[\alpha\,
   \frac{dt_1\wedge\cdots\wedge dt_b}
        {\prod_{i=1}^b(t_i-v_i)}\right],
 \qquad \alpha\in\omega_Y(R+D_h).
\end{equation}
A form regular in the $R$-support direction gives zero, even though
poles on $D_h$ are allowed. This computes the absolute dualizing
sheaf of the whole reduced divisor, including its crossings.

The map is independent of the lifts. If $v_i$ changes by a multiple
of $r_0$, the correction to the simple graph fraction becomes
regular in the support direction and vanishes. Under a base
coordinate change the Jacobian in the top differential cancels
the determinant in the ordered graph residues. These are also the
usual transformation rules for a complete-intersection residue.
Thus~\eqref{bl:hod:simple-graph} is intrinsic.

For completeness, strictness applies here to a projective morphism,
despite the quasi-projective ambient notation. Compactify the
projection projectively over $B$. Since $Q\to B$ is proper, its
image in that compactification remains closed and acquires no
points in the relative boundary. Extend the object with this
unchanged support. Projective filtered strictness applies to it
and hence computes the original direct image. The mixed statement
is Saito's Theorem~2.14 and Proposition~2.15
in~\cite{Saito1990}, with the algebraic formulation in its
Section~4.1; no decomposition of a mixed object is needed.

At level zero,~\eqref{bl:hod:relative-dr} has only the last term.
Its first cohomological direct image is therefore
$R^1h_*\omega_R(D_h)$. Strictness identifies this with $F_0\mathcal N$
as a subsheaf of $\mathcal N$, and gives the assertion at negative
levels as well. Finally, base differentiation commutes with the
relative differential in the product. It therefore commutes with
the map from the last term to differential-module direct image.
\end{proof}

The same calculation describes every filtration level locally:
$F_j\mathcal M$ consists of the graph fractions with simple initial
denominators in the boundary and graph coordinates, and at most
$j$ extra denominator exponents in total, modulo the fractions
regular in the $R$-support direction. We will use this explicit
description across a ramified base change.

\subsection{The consequence of Kodaira--Saito vanishing}

For a mixed Hodge module $\mathcal N_0$ on a smooth projective variety
$T_0$ and an ample line bundle $H_0$, Kodaira--Saito vanishing gives
\begin{equation}\label{bl:hod:ksv}
 \mathbb H^i\!\left(T_0,
     \gr_j^F\DR(\mathcal N_0)\otimes H_0^{-1}\right)=0
 \qquad(i<0).
\end{equation}
We use perverse normalization, so the right de Rham complex ends
in degree zero. The mixed version follows from the pure strict-support decomposition
and the strict finite weight filtration; see~\cite[Proposition~2.33]{Saito1990}
and~\cite[Theorem~28]{Popa2016}.
Tate twists do not affect the assertion, which holds for every $j$.

\begin{lemma}[No ample line in the lowest symbol kernel]
\label{bl:lem:symbol-vanishing}
Suppose $F_{<0}\mathcal N_0=0$. Then
\begin{equation}\label{bl:hod:symbol-vanishing}
 \Hom\!\left(H_0,
  \ker\left[F_0\mathcal N_0\otimes T_{T_0}
                \longrightarrow\gr_1^F\mathcal N_0\right]\right)=0.
\end{equation}
The arrow is the symbol of the right differential-operator action.
\end{lemma}

\begin{proof}
At $j=1$ all terms in degrees at most $-2$ vanish, so the graded
de Rham complex is the two-term complex
\[
 [F_0\mathcal N_0\otimes T_{T_0}\longrightarrow\gr_1^F\mathcal N_0]
 \quad\text{in degrees }-1,0.
\]
Its hypercohomology in degree $-1$, after tensoring by $H_0^{-1}$,
is exactly the space of global sections of the twisted kernel.
Equation~\eqref{bl:hod:ksv} proves the claim.
\end{proof}

\begin{remark}\label{bl:hod:whole-sheaf}
Both Lemma~\ref{bl:lem:lowest-piece} and
Lemma~\ref{bl:lem:symbol-vanishing} concern complete coherent sheaves,
including sections on smaller supports. They require no local
freeness and no restriction to the open set carrying a variation
of Hodge structure. This is essential for the obstruction calculation.
\end{remark}
